\[
y(t)=x(2)+x(t-6)
\]

\textbf{خطی بودن:}

برای ورودی‌های دلخواه \(x_1(t)\) و \(x_2(t)\) و ضرایب \(a,b\) داریم:
\[
T\{a x_1+b x_2\}(t)
=(a x_1+b x_2)(2)+(a x_1+b x_2)(t-6)
\]
\[
=a\bigl(x_1(2)+x_1(t-6)\bigr)+b\bigl(x_2(2)+x_2(t-6)\bigr)
=a\,T\{x_1\}(t)+b\,T\{x_2\}(t).
\]
پس سیستم {خطی} است.

$\\$
\textbf{علّیت:}

برای علّی بودن، \(y(t)\) باید فقط به مقادیر \(x(\tau)\) با \(\tau\le t\) وابسته باشد.
برای مثال در \(t=1\):
\[
y(1)=x(2)+x(-5).
\]
جمله‌ی \(x(2)\) به زمان \(2>1\) وابسته است؛ بنابراین برای محاسبه‌ی \(y(1)\) به مقدار ورودی در آینده‌ی نسبی نیاز داریم. پس سیستم {غیرعلّی} است.

$\\$
\textbf{تغییرناپذیری زمانی:}

ورودی تأخیرخورده \(x_d(t)=x(t-t_0)\) را در نظر بگیرید:
\[
y_d(t)=x_d(2)+x_d(t-6)=x(2-t_0)+x(t-6-t_0).
\]
خروجی متناظر با تأخیر \(t_0\) روی \(y\) برابر \(y(t-t_0)=x(2)+x(t-t_0-6)\) است.
این دو عبارت برای \(t_0\neq 0\) (به‌طور کلی) برابر نیستند؛ پس سیستم {با زمان متغیر} است.

$\\$
\textbf{پایداری:}

اگر برای همه‌ی \(t\) داشته باشیم \(|x(t)|\le M\)، آنگاه \(|x(2)|\le M\) و \(|x(t-6)|\le M\)، و:
\[
|y(t)|\le |x(2)|+|x(t-6)|\le 2M.
\]
بنابراین هر ورودی کران‌دار، خروجی کران‌دار می‌دهد و سیستم {پایدار} است.
